\part{A Revised Algorithmic Workflow for Stage Stratification}

This section details a multi-phase, algorithmic workflow for implementing the data-driven curriculum. The process is designed to be systematic and reproducible, transforming the raw simulation outputs into a final, robust stage assignment for every trajectory in the dataset. The workflow is divided into four phases: (I) Invariant Computation and Feature Engineering, (II) Exploratory Data Analysis, (III) Derivation of Adaptive Stage Criteria, and (IV) Final Stage Assignment and Balancing.

\section{Phase I: Invariant Computation and Feature Engineering}

The foundation of the data-driven curriculum is a comprehensive and accurate feature set. This phase focuses on processing the raw simulation data to compute the physical invariants discussed in Part 1 and correcting for any known sampling artifacts.

\subsection{Pre-computation and Data Loading}

This initial step leverages the project's existing data processing pipeline. It involves executing the established Python scripts to perform the following actions for every simulation in the dataset:
\begin{enumerate}
    \item \textbf{Parse Raw Data:} Read the initial condition files (\texttt{ic\_*.txt}) and the corresponding Brutus N-body integrator output files (\texttt{brutus\_*.txt}).[5]
    \item \textbf{Center-of-Mass (COM) Transformation:} Verify and, if necessary, transform the initial conditions to the COM frame, ensuring zero total momentum. This is a critical prerequisite for consistent physical calculations, especially since the `hierarchical` ICs may not initially be in the COM frame.[5]
    \item \textbf{Scaling to Dimensionless Units:} Apply the global characteristic scales for length, time, and velocity to convert all physical quantities into a consistent, dimensionless system where $G=1$. This is essential for numerical stability and for ensuring the PINN inputs are on a comparable scale.[5]
\end{enumerate}
The output of this step is a clean, standardized set of initial conditions (masses, positions, velocities) for every simulation.

\subsection{Invariant Calculation}

For each standardized initial condition, a dedicated function should compute the full suite of physical invariants defined in Section 1.2. This includes:
Total Mechanical Energy ($E$), Raw Virial Ratio ($Q_{\text{raw}}$), Normalized Angular Momentum ($L_{\text{norm}}$), Minimum Initial Separation ($s_{\text{min}}$), Shape Invariants ($I_1, I_2, I_3$), and Lyapunov Time ($\tau_L$), estimated from the corresponding Brutus trajectory data using the Rosenstein-style estimator detailed in the project documentation.[5]

\subsection{Correcting for Sampling Artifacts: The Normalized Virial Ratio}

A critical issue identified in the project is the distortion of the virial ratio for quasi-random ICs due to the Sobol sampling scale factor, $s$.[5] To address this, a normalized virial ratio, $Q_{\text{norm}}$, must be computed to place all simulations on a comparable footing.

The transformation is based on the physical scaling of the kinetic and potential energies. Since kinetic energy $T \propto v^2$ and potential energy $V \propto 1/r$, and the sampling scales positions as $r \to sr$ and velocities as $v \to sv$, the virial ratio scales as:
$$Q \propto \frac{T}{|V|} \propto \frac{(sv)^2}{1/(sr)} = s^3 \frac{v^2}{1/r}$$
Therefore, to recover a scale-invariant measure of virial equilibrium, the raw virial ratio must be divided by $s^3$. The normalized virial ratio is defined as:
$$Q_{\text{norm}} = \frac{Q_{\text{raw}}}{s^3}$$
This calculation should be applied to all simulations of type `quasi-random`. For `hierarchical` simulations, it can be assumed that $s=1$, so $Q_{\text{norm}} = Q_{\text{raw}}$.

\authcomment{I'm not sure what this s is referring to. Should we look for the exact value used in the sampling script?}

\textbf{Contingency Plan:} If the scale factor $s$ is not explicitly stored with each IC, it must be recorded during the IC generation or estimated post-hoc. As suggested in the project analysis, one method is to estimate it by comparing the median momentum of the scaled sample to a baseline unscaled sample.[5] Without this correction, the virial ratio is not a reliable metric for the quasi-random set and should be excluded from the analysis for those simulations.

\subsection{Assembling the Invariant Feature DataFrame}

The final step of this phase is to consolidate all computed information into a single, master \texttt{pandas.DataFrame}. This DataFrame will be the central data structure for all subsequent analysis and serves as a persistent, intermediate representation of the dataset's physical properties.

\begin{table}[H]
\centering
\caption{Recommended structure for the Invariant Feature DataFrame.}
\label{tab:dataframe}
\resizebox{\textwidth}{!}{%
\begin{tabular}{@{}lllllllllll@{}}
\toprule
\textbf{sim\_id} & \textbf{ic\_type} & \textbf{s\_factor} & \textbf{E} & \textbf{Q\_raw} & \textbf{Q\_norm} & \textbf{L\_norm} & \textbf{s\_min} & \textbf{tau\_L} & \textbf{I1} & \textbf{I2} \\ \midrule
0001 & hierarchical & 1.0 & -2.5 & 1.01 & 1.01 & 0.95 & 3.1 & 1000.0 & 0.8 & 0.9 \\
5001 & quasi-random & 1.2 & -0.8 & 1.75 & 1.01 & 0.45 & 1.5 & 12.3 & 0.3 & 0.7 \\ \bottomrule
\end{tabular}%
}
\end{table}

This DataFrame provides a complete, queryable summary of the entire dataset, enabling the powerful exploratory and algorithmic techniques described in the following phases.

\section{Phase II: Exploratory Data Analysis (EDA) and Landscape Mapping}

Before defining the final curriculum stages, it is essential to perform a thorough exploratory data analysis (EDA). The goal of this phase is to build a strong intuition for the structure of the data, validate the physical assumptions, and identify the key relationships that will guide the stratification strategy. This process transforms the numerical data in the invariant DataFrame into actionable visual insights.

\subsection{Stratified Distribution Analysis}

The first step is to understand the distribution of each key difficulty metric. This is achieved by generating histograms and, for a smoother representation, Kernel Density Estimate (KDE) plots for the primary invariants: $E$, $Q_{\text{norm}}$, $L_{\text{norm}}$, and $\tau_L$.

Crucially, these visualizations must be \textbf{stratified by \texttt{ic\_type}}. Plotting the distributions for `hierarchical` and `quasi-random` systems separately (or overlaid with different colors) is expected to reveal fundamental differences between the two datasets.[5]

\subsection{Correlation Analysis}

Next, the relationships \textit{between} the different difficulty metrics must be investigated. A Pearson correlation matrix, visualized as a heatmap, is the standard tool for this analysis. Again, this should be performed separately for the `hierarchical` and `quasi-random` subsets.

\subsection{Low-Dimensional Visualization}

To gain a more holistic view of the difficulty landscape, it is useful to visualize the data in a low-dimensional space.
\begin{itemize}
    \item \textbf{Pair Plots:} A \texttt{seaborn.pairplot} can be used to show all pairwise scatter plots between the key invariants, with points colored by \texttt{ic\_type}.
    \item \textbf{PCA Projection:} A powerful visualization technique is to perform a preliminary PCA on the standardized invariants and create a 2D scatter plot of the first two principal components (PC1 vs. PC2). The points on this plot can be colored by a key chaos indicator, such as $\tau_L$ or its inverse, $\lambda_{\text{max}}$.
\end{itemize}
This EDA phase is not merely a formality; it is a crucial step that ensures the subsequent algorithmic choices are grounded in the empirical reality of the data, not just in theoretical assumptions.

\section{Phase III: Deriving Adaptive Stage Criteria - A Hybrid Strategy}

With a deep understanding of the data's structure from the EDA phase, the next step is to derive the specific rules for stage assignment. The guiding principle of the revised framework is to use a hybrid strategy tailored to the distinct physical nature of the two sub-datasets.

\begin{table}[H]
\centering
\caption{Comparison of Original vs. Proposed Data-Driven Curriculum Framework}
\label{tab:framework_comparison}
\begin{tabular}{@{}p{2.5cm}p{5.5cm}p{5.5cm}@{}}
\toprule
\textbf{Aspect} & \textbf{Original Fixed-Threshold Framework} & \textbf{Proposed Data-Driven Framework} \\ \midrule
\textbf{Thresholds} & Static, hard-coded values based on theory (e.g., $Q<1.3$, $\tau_L>5$).[5] & Adaptive, derived from data statistics (e.g., quantiles, cluster centroids).[5] \\
\textbf{Stage Population} & Prone to severe imbalance or empty stages if data distribution mismatches theory.[5] & Balanced by design through quantile binning and population checks, ensuring smooth pacing.[5] \\
\textbf{Robustness} & Brittle; fails when data generation artifacts (e.g., scaling) violate assumptions.[5] & Robust; automatically adapts to the dataset's empirical distribution, including skews and artifacts.[5] \\
\textbf{Difficulty Ordering} & Assumes a simple, monotonic difficulty gradient that may not hold in multi-dimensional feature space. & Discovers natural difficulty gradients and qualitatively distinct regimes via PCA and clustering.[5, 13] \\
\textbf{Interpretability} & High, as criteria are based on physical theory. & High, as data-driven clusters and quantiles are interpreted and labeled by their physical properties (e.g., mean $\tau_L$). \\ \bottomrule
\end{tabular}
\end{table}

\subsection{Stratifying the Stable Subset (Hierarchical ICs): Stages 0 and 1}

The `hierarchical` dataset consists entirely of dynamically stable systems. The goal here is not to separate stable from chaotic, but to rank the systems by their \textit{degree} of stability and regularity. A simple, physically interpretable approach using quantile-based thresholds is most appropriate.

\subsection{Stratifying the Complex Subset (Quasi-Random ICs): Stages 2-6}

The `quasi-random` dataset is far more heterogeneous. A more powerful, multi-step strategy is required to properly parse this complexity. This involves a pipeline of feature preparation, dimensionality reduction with PCA, and regime discovery with unsupervised clustering. The algorithm involves:
\begin{enumerate}
    \item \textbf{Feature Preparation:} Select features $[E, Q_{\text{norm}}, L_{\text{norm}}, s_{\text{min}}, 1/\tau_L]$. Use the inverse Lyapunov time, $\lambda_{\text{max}} = 1/\tau_L$, to handle regular orbits where $\tau_L \to \infty$. Standardize all features.
    \item \textbf{Dimensionality Reduction with PCA:} Apply PCA to the standardized features to create a composite, low-dimensional representation.
    \item \textbf{Discovering Regimes with Clustering:} Apply K-Means (or GMM) to the principal component scores. The number of clusters, $k$, can be set to 5 (for Stages 2-6) or determined using heuristics like the elbow method.
    \item \textbf{Ordering Clusters by a Chaos Proxy:} Sort the discovered clusters by their median $\tau_L$ in descending order. The cluster with the highest median $\tau_L$ is the least chaotic and will form the earliest stage in this block.
\end{enumerate}

\section{Phase IV: Final Stage Assignment and Population Balancing}

The final phase translates the derived criteria into a definitive assignment of a stage label to every simulation.

\subsection{The Definitive Assignment Algorithm}

The final stage assignment rules are summarized in Table \ref{tab:final_stages}.

\subsection{Population Control and Stage Merging}

After the initial assignment, it is imperative to check the population of each stage. If any stage falls below a minimum population threshold (e.g., 5\% of the total dataset), it should be merged with its most similar adjacent stage to ensure the curriculum has a sensible "pacing".

\begin{longtable}{p{1.2cm}p{2.8cm}p{3.2cm}p{3.8cm}p{4.5cm}}
\caption{Final Proposed Stage Definitions and Data-Driven Criteria}
\label{tab:final_stages} \\
\toprule
\textbf{Stage} & \textbf{Primary Data Source} & \textbf{Stratification Method} & \textbf{Specific Data-Driven Criteria} & \textbf{Expected Physical Characteristics} \\
\midrule
\endfirsthead

\multicolumn{5}{c}%
{{\tablename\ \thetable{} -- continued from previous page}} \\
\toprule
\textbf{Stage} & \textbf{Primary Data Source} & \textbf{Stratification Method} & \textbf{Specific Data-Driven Criteria} & \textbf{Expected Physical Characteristics} \\
\midrule
\endhead

\midrule
\multicolumn{5}{r}{{Continued on next page}} \\
\endfoot

\bottomrule
\endlastfoot

\textbf{0} & Hierarchical & Quantile Binning & $E$ in bottom quartile \& $L_{\text{norm}}$ in top quartile. & Deeply bound, near-circular, highly regular orbits. \\[0.5em]

\textbf{1} & Hierarchical & Remainder & All other hierarchical systems. & Stable, bound, regular orbits, but less constrained. \\[0.5em]

\textbf{2} & Quasi-Random & PCA + Clustering & Members of the cluster with the highest median $\tau_L$. & Most stable of the chaotic set; long-term regular motion. \\[0.5em]

\textbf{3} & Quasi-Random & PCA + Clustering & Members of the cluster with the 2nd highest median $\tau_L$. & Onset of chaos; mildly chaotic systems. \\[0.5em]

\textbf{4} & Quasi-Random & PCA + Clustering & Members of the cluster with the 3rd highest median $\tau_L$. & Moderately chaotic systems, possibly resonant. \\[0.5em]

\textbf{5} & Quasi-Random & PCA + Clustering & Members of the cluster with the 4th highest median $\tau_L$. & Strongly chaotic systems with short predictability. \\[0.5em]

\textbf{6} & Quasi-Random & PCA + Clustering + Outlier & Most chaotic cluster (lowest median $\tau_L$) + outliers. & Violent, unpredictable dynamics; ejections, unbound systems. \\[0.5em]

\end{longtable}
